% Expository recovery from the source owners of the TPK kernel.
% Provenance and preservation checks are documented in technical/TPK_INTEGRACION.md.
% This file extends the TPK section; it does not replace extension.tex or generacion.tex.

\subsubsection{Dynamical state, observables, and dependence of operations}
\label{tpk-int:estado}

The topological phase kernel organizes the evolution of the two APP sheets
with the orientation determined by TRIT. Its definition comprises the state
on which it acts, the transition rules, and the relations making
prolongations compatible. Equation \eqref{eq:actualizacion} acquires
operational content when the coordinates received and modified by each
of its factors are specified. Position belongs to the toroidal support; phase
belongs to the schedule; orientation determines transport; quotients
record the arithmetic evaluation before reduction.
These coordinates are related through the transition and preserve different
mathematical functions.

The observable projection of a state is
\begin{equation}
 q=(x_+,d_+,x_\times,d_\times,\theta)
 \in\mathcal Q_{\rm obs}
 :=(X_9\times\mathcal D)^2\times\Theta_{108},
 \qquad \mathcal D=\{N,E,S,O\},\quad
 \Theta_{108}=\ZZ/108\ZZ.
 \label{tpk-int:qobs}
\end{equation}
The subscripts distinguish the additive and multiplicative cursors. The
fibres of complete states over \(q\) preserve the ordered trajectory, the integer
evaluations of both sheets, their remainders and quotients, and the orientation
and boundary records. Precision prolongation adds the
prefix and compatible cylinder. Two states with the same
projection \(q\) can therefore retain different records.

The dependence of operations has three levels. At the local level,
the selector activates a sheet, transport modifies its position, and
updating calculates the new arithmetic coordinates. At the trajectory
level, these transitions compose and preserve the records of
preceding steps. At the refinement level, extension relations
determine which continuation preserves the prefix, signature, and
boundary of its antecedent. The finite schedule describes the phase of these
operations; holonomy describes their composition on the history.

\subsubsection{Tritic selection and lifted cardinal transport}
\label{tpk-int:seleccion}

Let \(\bar\theta\in\{0,\ldots,107\}\) be the representative of the schedule.
The operating phase and dodecaphasic sector are
\begin{equation}
 r(\theta)=1+(\bar\theta\bmod9),\qquad
 \beta(\theta)=\left[\left\lfloor\frac{\bar\theta}{9}\right\rfloor\right]_{12}.
 \label{tpk-int:fase}
\end{equation}
The selector already defined can be written as \(\varepsilon_9=M_{\rm ph}\).
Its value vector, in the ordered chart of \(D_9\), is
\[
 m_{\rm ph}=(1,-1,0,1,-1,0,1,-1,0).
\]
The function \(M_{\rm ph}:D_9\to\TRIT\) and this vector are two
presentations of the same selector, with their respective domains.

Define the activation indicators
\[
 a_+(\theta)=\mathbf1_{\{1,4,7\}}(r(\theta)),\qquad
 a_\times(\theta)=\mathbf1_{\{2,5,8\}}(r(\theta)).
\]
Cursor transport is then
\begin{equation}
 x_+'=x_++a_+(\theta)\nu_{d_+},\qquad
 x_\times'=x_\times+a_\times(\theta)\nu_{d_\times}
 \quad\text{in }X_9.
 \label{tpk-int:cursores}
\end{equation}
In phases \(3,6,9\), both positions remain unchanged; the neutral event
forms part of the chronology and its record. The cardinal directions
determine the vectors \(\nu_d\), whereas the phase signature determines
which of these vectors is applied.

The integer lift preserves boundary crossing. If
\(s:X_9\to D_9^2\) chooses the positive representatives and
\(x'=x+a\nu_d\), with \(a\in\{0,1\}\), define
\begin{equation}
 b_d(x,a)=\frac{s(x)+a\nu_d-s(x')}{9}\in\ZZ^2.
 \label{tpk-int:borde}
\end{equation}
Divisibility by \(9\) follows from equality in \(X_9\).
If \(z\in\ZZ^2\) records preceding crossings, the update is
\(z'=z+b_d(x,a)\). Thus,
\[
 s(x')+9z'=s(x)+9z+a\nu_d.
\]
The integer transport datum remains reconstructible after position is
reduced on the torus.

\begin{proposition}[Conservation of lifted displacement]
\label{tpk-int:prop-desplazamiento}
For a cardinal trajectory of \(n\) steps with activation indicators
\(a_t\),
\[
 s(x_n)+9z_n=s(x_0)+9z_0+
 \sum_{t=0}^{n-1}a_t\nu_{d_t}.
\]
The same equality is preserved when two compatible trajectories are concatenated.
\end{proposition}
\begin{proof}
The one-step identity is \eqref{tpk-int:borde}. Summing it over the
\(n\) steps cancels the intermediate coordinates. Splitting the summation
interval into two parts produces exactly the concatenation law. For a
closed trajectory, the difference \(z_n-z_0\) is the winding number
defined previously.
\end{proof}

The same principle governs APP evaluations. If a sheet takes the
integer value \(N_t\), its record simultaneously preserves
\(N_t=R_t+9q_t\). The difference between two steps satisfies
\[
 N_{t+1}-N_t=(R_{t+1}-R_t)+9(q_{t+1}-q_t).
\]
Spatial crossing memory \(z_t\) and the arithmetic quotient \(q_t\)
have different definitions; their joint preservation retains the geometric
and arithmetic provenance of each transition.

\subsubsection{Two-cursor chronology and composite returns}
\label{tpk-int:cronologia}

The cardinal involution \(\jmath\) exchanges \(N\) with \(S\) and
\(E\) with \(O\). After \eqref{tpk-int:cursores} is applied, both
directions are reversed at the end of a block of \(27\) steps.
Writing
\[
 h(\theta)=\mathbf1_{\{0,27,54,81\}}
             ((\bar\theta+1)\bmod108),
\]
the complete rule on the observable projection is
\begin{equation}
 F_{\rm obs}(q)=
 \bigl(x_+',\jmath^{h(\theta)}d_+,
       x_\times',\jmath^{h(\theta)}d_\times,
       \theta+1\bigr).
 \label{tpk-int:fobs}
\end{equation}
The readout forming the emission uses positions before
transport, according to the explicit recurrence of
Section~\ref{sec:extension}. Observable updating and emission
therefore compose in a fixed order.

\begin{proposition}[Reversibility and period of the observable schedule]
\label{tpk-int:periodo}
The map \(F_{\rm obs}\) is a permutation of
\(\mathcal Q_{\rm obs}\) and satisfies \(F_{\rm obs}^{108}=I\).
A block of \(54\) steps restores positions and directions;
the coordinate in \(\Theta_{108}\) advances by \(54\).
\end{proposition}
\begin{proof}
Given the final state, first recover the preceding phase by subtracting one
in \(\Theta_{108}\). That phase determines whether the
directions were reversed and which cursor moved. Applying the corresponding involution
and subtracting the cardinal vector recovers the initial state uniquely.

Each interval of \(27\) steps contains nine activations of each
cursor. The following block preserves the activations and reverses the
directions, so the integer displacements cancel over
\(54\) steps. The rule has period \(54\) in its position and
orientation increments; the same cancellation holds for any origin
of the schedule. After \(108\) steps, \(\theta\) is also restored.
\end{proof}

We reserve
\[
 \mathcal K_{27}=F_{\rm obs}^{27}:
 \mathcal Q_{\rm obs}\longrightarrow\mathcal Q_{\rm obs}
\]
for the iterate of \(27\) transitions. Its fourth iterate is the
identity on the observable state defined in \eqref{tpk-int:qobs}.
Projection of the history onto the schedule permits this finite
check. The trajectory record, its evaluations, and
new prefixes are preserved in the dynamical state on which
\(\mathcal U_t\) acts; observable equality does not reinitialize them.

\subsubsection{Affine memory transport and composition of states}
\label{tpk-int:memoria}

The observable projection admits fibres of arithmetic and geometric data.
Over a configuration \(q\), write a state in the form
\[
 x=(q,o,h,c,\sigma,C,b,\lambda).
\]
Here \(o\) preserves orientation and sheet; \(h\), remainders
and quotients; \(c\), additive transport memory; \(\sigma\),
the boundary signature; \(C\), the precision cylinder; \(b\), its
boundary label; and \(\lambda\), the recorded trajectory. The
signature is determined by a map
\[
 \operatorname{Sig}_q:
 \mathsf H_q\times\mathsf C_q\times\mathsf{Hist}_q
 \longrightarrow\mathsf S_q,
 \qquad \sigma=\operatorname{Sig}_q(h,c,\lambda).
\]
The symbols \(\mathsf H_q,\mathsf C_q,\mathsf{Hist}_q\) and
\(\mathsf S_q\) denote, respectively, the remainder--quotient data sets,
the abelian group of transport memory, the
recorded trajectories, and the boundary signatures.
Realizations of this signature are specified componentwise in
the regional and terminal constructions. This dependence prevents
the signature from being treated as a free parameter once the data producing it are fixed.

The elementary remainder--quotient chart is explicit. For
\(n\in\ZZ\), let
\[
 r=1+((n-1)\bmod9),\qquad u=\frac{n-r}{9}.
\]
Then \(n=r+9u\), with \(r\in D_9\) and \(u\in\ZZ\).
Uniqueness follows because two representatives of \(D_9\) congruent
modulo \(9\) are equal. This chart provides local integer
reconstruction. The precision towers additionally incorporate truncation
maps and the cylinders developed in
\ref{sec:extension} and \ref{sec:generacion}.

Let \(\gamma:q\to q'\) be an observable trajectory. Fibre
transport is expressed through maps
\[
 \mathsf H(\gamma):\mathsf H_q\to\mathsf H_{q'},\qquad
 \mathsf C(\gamma):\mathsf C_q\to\mathsf C_{q'},\qquad
 \mathsf{Hist}(\gamma):\mathsf{Hist}_q\to\mathsf{Hist}_{q'}.
\]
These maps respect identities and composition, and
\(\mathsf C(\gamma)\) is a homomorphism of abelian groups.
The increment \(k_{\rm tr}(\gamma)\in\mathsf C_{q'}\) satisfies
\begin{equation}
 \begin{split}
 k_{\rm tr}(\operatorname{id}_q)&=0,\\
 k_{\rm tr}(\gamma_2\circ\gamma_1)
 &=k_{\rm tr}(\gamma_2)+
   \mathsf C(\gamma_2)k_{\rm tr}(\gamma_1).
 \end{split}
 \label{tpk-int:cociclo}
\end{equation}
The symbol \(k_{\rm tr}\) here denotes the transport cocycle;
the rational coordinate \(\kappa\) of record \(K\) retains
its notation in the chapter on dodecaphasic closure.

The transportable coordinates update by
\begin{equation}
 \begin{aligned}
 h'&=\mathsf H(\gamma)h,&
 c'&=\mathsf C(\gamma)c+k_{\rm tr}(\gamma),\\
 \lambda'&=\gamma\circ\lambda,&
 \sigma'&=\operatorname{Sig}_{q'}(h',c',\lambda').
 \end{aligned}
 \label{tpk-int:accion-memoria}
\end{equation}
The second equation is an affine action: the linear term transports
the preceding memory and the cocycle adds the increment produced by the
trajectory. In the spatial-crossing realization for each cursor separately,
\(\mathsf C(\gamma)=I\) and \(k_{\rm tr}(\gamma)\) is the sum
of the vectors in \eqref{tpk-int:borde}. This concrete instance
links the abstract law to the cardinal operations already defined.

\begin{proposition}[Compatibility of updating with concatenation]
\label{tpk-int:composicion-memoria}
Updating according to \eqref{tpk-int:accion-memoria} by
\(\gamma_1\), followed by updating by \(\gamma_2\),
agrees with updating by \(\gamma_2\circ\gamma_1\).
\end{proposition}
\begin{proof}
For the affine coordinate, two updates give
\[
 c''=\mathsf C(\gamma_2)\mathsf C(\gamma_1)c
       +\mathsf C(\gamma_2)k_{\rm tr}(\gamma_1)
       +k_{\rm tr}(\gamma_2).
\]
Functoriality of \(\mathsf C\) and
\eqref{tpk-int:cociclo} turn this expression into
\(\mathsf C(\gamma_2\circ\gamma_1)c+
k_{\rm tr}(\gamma_2\circ\gamma_1)\).
The assertion for \(h\) follows from composition of
\(\mathsf H\), and that for \(\lambda\) from associativity of
trajectories. The final signature agrees because it is evaluated on the same
three reconstructed coordinates.
\end{proof}

The cylinder and its boundary complete this law through the relations
\(\operatorname{Ext}_\gamma\), defined in
\ref{tpk-int:bidireccionalidad}. Their elements determine which of
the preceding updates preserve an admissible prolongation.
The states, together with the triples \((\gamma,x,x')\) admitted by
these relations, form a category over observable trajectories:
\[
 (\gamma_2,x',x'')\circ(\gamma_1,x,x')
     =(\gamma_2\circ\gamma_1,x,x'').
\]
The identity of \(x\) is \((\operatorname{id}_q,x,x)\).
Composition preserves arithmetic transport and boundary
prolongability simultaneously. This is the structure in which returns
with memory and successive refinements take place.

\subsubsection{Dodecaphasic record, phase, and accumulated memory}
\label{tpk-int:dodecafase}

The \(108\) transitions are distributed among the ordered windows
\[
 E_m=\{9m-8,\ldots,9m\},\qquad 1\leq m\leq12.
\]
The record preserves the phase origin, orientation, and data
produced during each window. Its map is
\[
 \mathcal R_{12}(x)=
 \bigl((E_m,\tau_m,\sigma_m,\kappa_m,\Lambda_m)\bigr)_{m=1}^{12},
\]
where \(\tau_m\) is the window subroute, \(\sigma_m\) its
orientation, \(\kappa_m\) the integer boundary increment, and
\(\Lambda_m\) the local incidence record. We adopt the chart
developed in \ref{subsec:k-registro-integral}; its window indices
correspond to \(m=\beta+1\).
The domain comprises compatible TPK histories,
and the codomain their oriented temporal records.
Composition of segments and their increments is governed by
\eqref{tpk-int:accion-memoria}. The group \(C_{12}\), by
contrast, describes translation of the window index; \(\Theta_{108}\) preserves
the step schedule. This distinction specifies what transforms
and what is observed upon completing a cycle.

In coordinates \(\theta=9b+r\), with \(0\leq r<9\), an advance
of \(9\) steps preserves \(r\) and increases \(b\) by one.
Phase addition introduces the quotient
\[
 c_0(r,r')=\left\lfloor\frac{r+r'}9\right\rfloor,
\]
whose cocycle identity ensures associativity of composition.
The finite schedule preserves \([b]_{12}\); a revolution record
preserves \(b\in\ZZ\), or \(b\in\mathbb N_0\) for
a history with fixed origin. Section~\ref{sec:extension} proves
the nonsplit exact sequence \eqref{eq:extension108}. Its function is
to describe the relation between these coordinates before they are used to
construct the signed state.

The signed record receives phase and memory data through the
composition developed in Section~\ref{sec:k-registro}:
\begin{equation}
 R_{\rm sgn}=\Pi_H\circ\mathcal E_{108}^{90,120}
                         \circ\mathcal R_{12}.
 \label{tpk-int:registro-causal}
\end{equation}
Its domain is the terminal state with memory; its output belongs to
\(\ZZ^{12}\). The reversible transformation between this signed state
and \(K\), together with the rational readout of \(K\), is proved in
the corresponding chapter. Recording operations precede
these transformations. Reconstruction of a record from its cyclic
differences verifies its recoverability; its dynamical provenance
must be preserved through the events that produced it.

For the coded trace \((d_t)\) preserved by the record, the
event reader counts the codes \(\{2,4,6,8\}\) in each
window with orientation
\((+,+,+,-,-,-,+,+,+,-,-,-)\). These codes and the cardinal
symbols of \(\mathcal D\) are distinct coordinates of the
record. Chapter~\ref{sec:k-registro} develops this reader,
the integral decomposition \(1+5+6\), the congruences
characterizing its image, and its exact inverse. The inversion proof
establishes which data this chart preserves; the definition of
\(\mathcal R_{12}\) specifies the oriented history on which
it acts.

\subsubsection{Phase incidence and internal construction of the channels}
\label{tpk-int:incidencia}

The selector determines three subsets of \(D_9\):
\[
 H_+=\{1,4,7\},\qquad H_-=\{2,5,8\},\qquad H_0=\{3,6,9\}.
\]
Let \(H_{\rm act}=H_+\sqcup H_-\) be the set of active phases and
let \(O_{\rm ph}=D_9\setminus\{9\}\) be the chart preceding the marked
return. Denote by \(P_2(H_{\rm act})\) the set of unordered pairs
of active phases. Internal incidence is given by
\begin{equation}
 \mathcal F^{\rm ph}_6=H_{\rm act}\times P_2(H_{\rm act}),\qquad
 \mathcal F^{\rm ph}_8=O_{\rm ph}\times P_2(H_{\rm act}),\qquad
 \Theta_{\rm ph}=D_9\times P_2(H_{\rm act}).
 \label{tpk-int:fases90120}
\end{equation}
The indices describe the number of positions in the first component.
The construction uses the phases and the schedule origin already fixed;
subsequent exceptional realizations will transport these
incidences into their own domains.

\begin{proposition}[Cardinalities and oriented phase incidence]
\label{tpk-int:cardinales}
The sets in \eqref{tpk-int:fases90120} have cardinalities
\(90,120,135\), respectively. If
\[
 \partial_{\rm ph}=\Theta_{\rm ph}\times\{-1,+1\},\qquad
 E_{\rm ph}=\partial_{\rm ph}\sqcup\{\infty\},\qquad
 L_{\rm ph}=H_0\times P_2(H_{\rm act}),
\]
then
\[
 |\partial_{\rm ph}|=270,\quad |E_{\rm ph}|=271,\quad
 |L_{\rm ph}|=45,\quad
 |\partial_{\rm ph}\sqcup L_{\rm ph}|=315.
\]
Simultaneous translation of the phase origin and the marked return preserves
all these cardinalities and their inclusions.
\end{proposition}
\begin{proof}
There are six active phases, so
\(|P_2(H_{\rm act})|=\binom62=15\). Products by
\(6,8,9\) give \(90,120,135\). Orientation doubles \(135\);
the return mark adds one element; the three neutral phases give
\(3\cdot15=45\). A schedule translation transports each
factor by a bijection and carries the marked point to its image. Products
and disjoint unions thus preserve incidence.
\end{proof}

The cardinality \(271\) here has a realization through phase
incidence. The cubic completion of Section~\ref{sec:extension} has
its own construction of the shell \(1000-729\). The correspondence
between these realizations must preserve the maps and marks identifying
them, as well as cardinality. Similarly, the selector
\(\varepsilon_9\), the incidence \((90,120)\), and the extractor
\(\mathcal E_{108}^{90,120}\) preserve the domains and operations
just distinguished.

\subsubsection{Trajectory reversal and conjugate prolongations}
\label{tpk-int:bidireccionalidad}

Bidirectionality begins with an operation on trajectories before
scalar values are assigned to them. For
\(\gamma=(x_0;d_1,\ldots,d_n)\), with endpoint \(x_n\), define
\[
 \operatorname{rev}\gamma
   =(x_n;\jmath d_n,\ldots,\jmath d_1).
\]
The origin of the reversed trajectory is the endpoint of the original.
Composition satisfies
\[
 \operatorname{rev}^2\gamma=\gamma,\qquad
 \operatorname{rev}(\gamma_2\circ\gamma_1)
   =\operatorname{rev}\gamma_1\circ\operatorname{rev}\gamma_2,
\]
and the integer displacement changes sign. These identities follow
by reversing step order and replacing each cardinal vector
with its opposite. For a loop, this gives
\(\operatorname{wind}(\operatorname{rev}\gamma)
=-\operatorname{wind}(\gamma)\).

On records, reversal also requires transport of component order
and event orientation. Route reversal
and signature conjugation are distinguishable operations: their
composition is specified when applied to a concrete record. In
particular, the regional involution exchanging the closure
and propagation components preserves the self-scaling axis; its formula is
presented in Section~\ref{mono-ampl:regiones}. This regional operation
acts on a different domain from reversal of a cardinal trajectory.

Extensions over a route \(r:q\to q'\) are expressed through
relations between fibres of compatible states. If \(\mathcal F_q\)
denotes the fibre of complete states over \(q\), require
\[
 \operatorname{Ext}_r\subseteq\mathcal F_q\times\mathcal F_{q'},
 \qquad \Delta_{\mathcal F_q}\subseteq
        \operatorname{Ext}_{\operatorname{id}_q},
\]
\[
 \operatorname{Ext}_{r_2}\circ\operatorname{Ext}_{r_1}
       \subseteq\operatorname{Ext}_{r_2\circ r_1}.
\]
Their composition preserves the evaluations of both sheets, orientation,
quotients, prefix, and boundary. Route reversal does not
by itself imply that every extension relation is a bijection.
Conjugate prolongation is defined on the domain where transport of
these data satisfies the corresponding compatibilities.

\subsubsection{Nonadic holonomy and continuity of refinement}
\label{tpk-int:holonomia}

A phase revolution composes transitions on the state with memory:
\begin{equation}
 \Gamma_{9,k}=\mathcal U_{k+8}\circ\cdots\circ\mathcal U_k.
 \label{tpk-int:gamma}
\end{equation}
On an individualized history, this functional composition realizes
the holonomy relation; on the complete domain, we retain
\(\Gamma_{9,k}\subseteq\widehat X_k\times\widehat X_{k+9}\),
with the admissible prolongations developed below.
The cyclic base has nine phases, whereas
the fibre preserves trajectory, quotients, and boundary.
Monodromy is the action of one revolution of this base on the fibres;
holonomy \(\Gamma_{9,k}\) expresses its realization through the effective
TPK transitions.

If \(g\) observes phase and \(M\) the number of preserved revolutions,
\[
 g(\Gamma_{9,k}x)=g(x),\qquad M(\Gamma_{9,k}x)=M(x)+1.
\]
The first identity expresses phase closure and the second identifies
the new temporal depth. Coinductive prolongation additionally preserves
its cylinders and prefixes. Restrictions between depths
compose an inverse system; a compatible section brings together its
antecedents without replacing them with the latest observation. This
arithmetic of histories and its unit are developed in
\ref{hist:seccion}--\ref{hist:doble-varianza}; the regions and their
readers are constructed in Section~\ref{sec:generacion}.

Phase transfer on already constructed emissions has a different
domain. Let \(\mathcal A_+\) and \(\mathcal A_\times\) be the
alphabets of additive and multiplicative signatures, of cardinalities \(18\)
and \(26\). Their product identifies the catalogue \(\mathcal U_6\).
For \(f:\mathcal A_+\times\mathcal A_\times\to\RR\),
the sheet averages are
\[
 (E_+f)(s,e)=\frac1{18}\sum_{s'\in\mathcal A_+}f(s',e),\qquad
 (E_\times f)(s,e)=\frac1{26}\sum_{e'\in\mathcal A_\times}f(s,e').
\]
These emissions and their cardinalities are constructed through
the recurrence in Section~\ref{sec:extension}. Once they have been obtained,
the transfer operator is defined on
\(\mathcal U_6\times C_9\) by
\[
 P_+=E_+,\quad P_-=E_\times,\quad P_0=I,\qquad
 \mathcal K_{\rm ph}((u,r),(v,r+1))=P_{M_{\rm ph}(r)}(u,v),
\]
with zero weight for the remaining phase transitions. The index
\(r\) uses the positive representative of the cyclic class.
The phase order \((+1,-1,0)^3\) then produces the renewal
\[
 \mathcal K_{\rm ph}^{9}\big|_r=E_+E_\times.
\]
Indeed, both averages are idempotent, commute, and average different
factors; their composition assigns weight \(1/(18\cdot26)=1/468\)
to each emission. This equality characterizes a subsequent transfer
observable. History transport remains given by
\eqref{tpk-int:gamma}, with its memory and boundary.

The resulting hierarchy establishes the function of the following chapters.
Finite emission constructs the domain and regions; extension relations
preserve their prefixes; the nonadic connection prolongs them; the
dodecaphasic record determines the data entering \(K\) and
the closure of \(\alpha\). Numerical realization and incidence
realization receive this preceding structure. This order permits
each consequence to be presented in its chapter while preserving, from the formal kernel,
the objects that make it possible.
